% !TEX root = ../Main.tex
\chapter{一元二次不等式（不含参）}

一元二次不等式形如 $ax^2+bx+c>0$（或 $\ge,<,\le$，$a\ne 0$）。求解的核心是\textbf{画图}：把二次函数 $y=ax^2+bx+c$ 的抛物线画出来，看它在 $x$ 轴上方还是下方——这就是数形结合。

\section{二次函数的三种表达式}

\state{二次函数表达式的选取}{
    解题时按需要在三种形式之间切换：
    \begin{itemize}
        \item 一般式：$y=ax^2+bx+c$（$a\ne 0$）；
        \item 顶点式：$y=a\pr{x+\dfrac{b}{2a}}^2+c-\dfrac{b^2}{4a}$；
        \item 两根式：$y=a(x-x_1)(x-x_2)$（前提 $\Delta\ge 0$，即有实根）。
    \end{itemize}
}

通常的做法：先把不等式整理成一般式，算判别式 $\Delta=b^2-4ac$，判断有没有实根。\textbf{若有根}（$\Delta\ge 0$）：

\begin{enumerate}
    \item 先尝试因式分解（写成两根式），再数形结合读解集；
    \item 分解不出来时，再用求根公式 $x=\dfrac{-b\pm\sqrt{\Delta}}{2a}$ 求根。
\end{enumerate}

\section{三个例子：按 $\Delta$ 的三种情形}

下面三题分别对应 $\Delta<0$、$\Delta>0$、$\Delta=0$，方法都是同一句话——画出抛物线，看 $x$ 轴上下。

\ex{解 $x^2-x+12>0$。}

\sol{
    $a=1>0$，开口向上；$\Delta=(-1)^2-4\cdot 1\cdot 12=-47<0$，无实根。抛物线整条都在 $x$ 轴上方，$y$ 恒正，所以不等式对一切 $x$ 都成立。

    解集 $\RR=(-\infty,+\infty)$。
}

\begin{center}
\begin{tikzpicture}[>=Stealth, scale=0.5]
    \draw[->] (-2.2,0) -- (3,0) node[right] {$x$};
    \draw[->] (0,-0.4) -- (0,2.9) node[above] {$y$};
    \draw[thick, blue, domain=-1.7:2.7, samples=50, smooth] plot (\x, {(\x-0.5)*(\x-0.5)+0.8});
    \node at (1.6, 2.1) [right] {\footnotesize $y=x^2-x+12$};
\end{tikzpicture}
\end{center}

\ex{解 $x^2-x-12>0$。}

\sol{
    $a=1>0$，开口向上；分解 $x^2-x-12=(x+3)(x-4)$，两根 $-3,4$。$y>0$ 即图象在 $x$ 轴上方，取两根的外侧。

    解集 $(-\infty,-3)\cup(4,+\infty)$。
}

\begin{center}
\begin{tikzpicture}[>=Stealth, scale=0.55]
    \draw[->] (-3.2,0) -- (3.2,0) node[right] {$x$};
    \draw[->] (0,-1.6) -- (0,1.8) node[above] {$y$};
    \draw[thick, blue, domain=-2.5:2.5, samples=60, smooth] plot (\x, {0.45*(\x+1.2)*(\x-1.8)});
    \filldraw (-1.2,0) circle (1.3pt) node[above left] {\footnotesize $-3$};
    \filldraw (1.8,0) circle (1.3pt) node[above right] {\footnotesize $4$};
\end{tikzpicture}
\end{center}

\ex{解 $x^2-14x+49>0$。}

\sol{
    $x^2-14x+49=(x-7)^2$，$\Delta=0$，重根 $x=7$。抛物线开口向上、与 $x$ 轴相切于 $x=7$，除切点外恒正。$x=7$ 时 $y=0$ 不满足严格不等式。

    解集 $\{x\mid x\ne 7\}=(-\infty,7)\cup(7,+\infty)$。
}

\begin{center}
\begin{tikzpicture}[>=Stealth, scale=0.5]
    \draw[->] (-0.8,0) -- (3.6,0) node[right] {$x$};
    \draw[->] (0,-0.3) -- (0,2.5) node[above] {$y$};
    \draw[thick, blue, domain=0.45:3, samples=50, smooth] plot (\x, {2.2*(\x-1.75)*(\x-1.75)});
    \filldraw (1.75,0) circle (1.3pt) node[below] {\footnotesize $7$};
\end{tikzpicture}
\end{center}
